\emph{Översikt}
I den här övningen använder du filhanteringen på riktiga uppgifter:
en fil full av text ska räknas och sammanställas, en fil ska översättas
till rövarspråket och tillbaka igen, och en databas med fortunekakor ska
läsas ur en fil i stället för att stå i programmet.  Varje uppgift har
ett lösningsförslag, och du får mest ut av materialet om du prövar
uppgiften själv först och sedan jämför ditt försök med förslaget.
Uppgifterna bygger på varandra: i den första samlar vi de frågor om
filer som alla de andra programmen behöver --- finns filen?  får den
skrivas över? --- i en egen modul, och sedan använder vi den modulen om
och om igen.
De tre sista uppgifterna är fördjupning och kan hoppas över.

\emph{Lärandemål}
Efter den här övningen ska studenten kunna:
% Lo08 (interaktiva program), Lo10 (data mellan fil och program)
\begin{restatable}{lo}{OvningFilesLOMissing}\label{OvningFilesLOMissing}%
  Hantera fel vid filhantering, som att filen inte finns, så att
  användaren får en begriplig fråga i stället för en krasch.
\end{restatable}
% Lo10 (data mellan fil och program)
\begin{restatable}{lo}{OvningFilesLOFileTypes}%
  \label{OvningFilesLOFileTypes}%
  Skilja mellan textfiler och binärfiler, och välja rätt läge när filen
  öppnas.
\end{restatable}
% Lo10 (data mellan fil och program)
\begin{restatable}{lo}{OvningFilesLOContextMgr}%
  \label{OvningFilesLOContextMgr}%
  Öppna och stänga filer på rätt sätt med \mintinline{python}{open} och
  \mintinline{python}{with}, och förklara varför filen måste stängas.
\end{restatable}
% Lo10 (data mellan fil och program), Lo09 (sammansatta datatyper)
\begin{restatable}{lo}{OvningFilesLORead}\label{OvningFilesLORead}%
  Läsa data från en fil, tolka det som står där och lagra det i lämpliga
  behållare.
\end{restatable}
% Lo10 (data mellan fil och program)
\begin{restatable}{lo}{OvningFilesLOWrite}\label{OvningFilesLOWrite}%
  Skriva data från programmet till en fil i ett format som programmet
  senare kan läsa tillbaka.
\end{restatable}
% Lo10 (data mellan fil och program), Lo06 (flexibla program)
\begin{restatable}{lo}{OvningFilesLOFormats}\label{OvningFilesLOFormats}%
  Läsa och skriva formatet CSV med hjälp av Pythons standardbibliotek och
  dess dokumentation.
\end{restatable}
% Lo10 (data mellan fil och program); övningsspecifik delaspekt av
% lärandemålet om filformat i föreläsningen Arbeta med filer.
\begin{restatable}{lo}{OvningFilesLOOwnFormat}%
  \label{OvningFilesLOOwnFormat}%
  Bestämma ett eget filformat för data som inte får plats på en rad, och
  skriva den kod som tolkar formatet.
\end{restatable}

\emph{Förkunskaper}
Studenten bör ha tagit del av föreläsningen \emph{Arbeta med filer}:
studenten ska kunna öppna och stänga filer med \mintinline{python}{open}
och \mintinline{python}{with}, läsa en fil både i sin helhet och rad för
rad, skriva till en fil, och veta att det som läses ur en textfil är
strängar.  Uppgifterna använder dessutom det som föreläsningarna
\emph{Inmatning och felhantering} (att fånga ett särfall med
\mintinline{python}{try} och \mintinline{python}{except} och att läsa en
spårutskrift), \emph{Villkor och styrstrukturer} (att fråga om igen när
svaret inte duger), \emph{Behållare: Listor},
\emph{Behållare: Uppslagslistor}, \emph{Upprepningar},
\emph{Funktioner} och \emph{Moduler och paket} (att lägga funktioner i en
egen modul och importera den) går igenom.
